\section{Kesimpulan}

Penelitian ini merancang dan menguji CoDA-EDA, yaitu EDA \textit{compact} yang mempertahankan dependensi antaratribut untuk \textit{ABAC policy mining} pada perangkat \textit{edge-IoT}.

Pertama, CoDA-EDA menggabungkan \textit{Probability Vector}, struktur dependensi Chow–Liu orde satu, serta pemisahan pembaruan parameter dan struktur. 
Rancangan ini memungkinkan pemodelan dependensi tanpa menyimpan populasi eksplisit.

Kedua, H1 didukung secara bersyarat. 
CoDA-EDA memberikan kualitas yang kompetitif dan menunjukkan keunggulan terhadap keluarga \textit{compact} univariat terutama pada ruang solusi berdimensi tinggi. 
Pada anggaran aturan yang memadai, kualitasnya sebanding dengan iBOA/BOA, tetapi masih dapat tertinggal dari ACO dan UBK pada penambangan statis.

Ketiga, H2 didukung kuat. 
Jejak memori CoDA-EDA tetap stabil terhadap ukuran populasi virtual dengan rasio $1,00$, sedangkan BOA, ACO, dan UBK meningkat $6,2$–$9,2$ kali ketika ukuran populasi dinaikkan enam belas kali.

Keempat, H3 dan H4 menunjukkan manfaat pada pembaruan kebijakan. 
Pembaruan struktur berbasis \textit{drift} mengurangi waktu komputasi ketika distribusi relatif stabil, meskipun manfaat H3 bergantung pada tingkat \textit{drift} dan lantai derau dataset. 
Sementara itu, \textit{warm-start} mempercepat re-optimasi sekitar $2,0$–$2,6$ kali dibandingkan \textit{re-mining} penuh dengan kualitas yang sekurang-kurangnya sebanding.

Kelima, validasi pada \textit{Smart Building IoT} menunjukkan pola hasil yang konsisten dengan pengujian utama, sehingga mendukung ketahanan temuan pada \textit{generator} sintetis yang berbeda, meskipun belum membuktikan generalisasi terhadap data operasional \textit{edge-IoT}.

Secara keseluruhan, kontribusi utama CoDA-EDA terletak pada efisiensi memori dan pembaruan kebijakan secara inkremental, bukan pada dominasi kualitas penambangan statis. 
Karakteristik tersebut menjadikannya sesuai untuk perangkat \textit{edge} dengan sumber daya terbatas yang memerlukan pembaruan kebijakan secara berkala.